\documentclass[12pt]{article}
\usepackage{amsmath,amssymb,amsthm,mathtools}
\usepackage[margin=1in]{geometry}
\usepackage{url}

\newtheorem{theorem}{Theorem}[section]
\newtheorem{lemma}[theorem]{Lemma}
\newtheorem{proposition}[theorem]{Proposition}
\newtheorem{corollary}[theorem]{Corollary}
\theoremstyle{definition}
\newtheorem{definition}[theorem]{Definition}
\theoremstyle{remark}
\newtheorem{remark}[theorem]{Remark}

\DeclareMathOperator{\supp}{supp}
\DeclareMathOperator{div}{div}
\newcommand{\D}{\mathcal D'}
\newcommand{\Cinf}{C^\infty}
\newcommand{\vol}{\operatorname{vol}}
\newcommand{\htop}{h_{\mathrm{top}}}

% The following web source is used only in the final remark concerning
% Pesin's entropy formula:
% [ar5iv.labs.arxiv.org](https://ar5iv.labs.arxiv.org/html/1002.2366)

\begin{document}

\title{Globally hypoelliptic vector fields have zero topological entropy}
\author{}
\date{}
\maketitle

\begin{abstract}
We prove that if a smooth vector field $X$ on a compact connected manifold
is globally hypoelliptic, in the sense that every distributional solution of
$\mathcal L_Xu=f$ with $f$ smooth is itself smooth, then the flow generated
by $X$ has zero topological entropy. The proof uses the Greenfield--Wallach
spectral theorem for globally hypoelliptic vector fields and the
zero-entropy theorem for systems with discrete spectrum.
\end{abstract}

\section{Conventions}

Throughout, $M$ is a compact connected smooth manifold, without boundary,
and $X\in C^\infty(M;TM)$. Its flow is denoted by
\[
g_t:M\longrightarrow M,\qquad \partial_t g_t(x)=X(g_t(x)).
\]

We fix an auxiliary Riemannian metric and denote its smooth positive
volume density by $d\vol$. Distributions are identified with generalized
densities. Thus a distribution $T$ is called smooth if there is a smooth
function $\rho$ such that
\[
\langle T,\varphi\rangle=\int_M\varphi\,\rho\,d\vol
\qquad(\varphi\in C^\infty(M)).
\]

The Lie derivative on distributions is defined by
\[
\langle\mathcal L_XT,\varphi\rangle
   :=-\langle T,X\varphi\rangle .
\]
The hypothesis of the problem is the following.

\begin{definition}[Global hypoellipticity]\label{def:GH}
The vector field $X$ is \emph{globally hypoelliptic} if
\[
u\in\D(M),\quad \mathcal L_Xu\in C^\infty(M)
\quad\Longrightarrow\quad u\in C^\infty(M).
\]
\end{definition}

For a flow, its topological entropy is by definition the topological
entropy of its time-one map:
\[
\htop(g):=\htop(g_1).
\]
This is the standard definition; see \cite[Chapter~3]{KatokHasselblatt}.

\section{Invariant measures are smooth and fully supported}

Every Borel probability measure $\mu$ on $M$ defines a distribution
$T_\mu$ by
\[
\langle T_\mu,\varphi\rangle=\int_M\varphi\,d\mu .
\]

\begin{lemma}\label{lem:invariant-measures-smooth}
If $\mu$ is invariant under the flow $g_t$, then $\mathcal L_XT_\mu=0$.
Consequently, under Definition~\ref{def:GH}, there is a smooth function
$\rho_\mu\geq 0$ such that
\[
d\mu=\rho_\mu\,d\vol.
\]
\end{lemma}

\begin{proof}
Flow invariance gives, for every $\varphi\in C^\infty(M)$ and every $t$,
\[
\langle T_\mu,\varphi\circ g_t\rangle
 =\int_M\varphi\circ g_t\,d\mu
 =\int_M\varphi\,d\mu
 =\langle T_\mu,\varphi\rangle .
\]
Differentiating at $t=0$ gives
\[
\langle T_\mu,X\varphi\rangle=0.
\]
Hence
\[
\langle\mathcal L_XT_\mu,\varphi\rangle
=-\langle T_\mu,X\varphi\rangle=0.
\]
Thus $\mathcal L_XT_\mu=0$. Since the right-hand side is smooth, global
hypoellipticity implies that $T_\mu$ is represented by a smooth function
$\rho_\mu$. Nonnegativity follows because $\mu$ is a positive measure.
\end{proof}

The stationary density satisfies the usual transport equation
\begin{equation}\label{eq:stationary-density}
X\rho_\mu+\rho_\mu\,\div_{\vol}X=0.
\end{equation}
Indeed,
\[
\mathcal L_X(\rho_\mu\,d\vol)
 =\bigl(X\rho_\mu+\rho_\mu\div_{\vol}X\bigr)d\vol.
\]

We shall need the following elementary topological observation.

\begin{lemma}\label{lem:boundary-nowhere}
If $U\subset M$ is open, then $\partial U$ has empty interior.
\end{lemma}

\begin{proof}
Suppose $V\subset\partial U$ is open and nonempty. Since
$V\subset\overline U$, we have $V\cap U\neq\varnothing$. On the other hand,
$V\subset\partial U$ and $U\cap\partial U=\varnothing$, so
$V\cap U=\varnothing$, a contradiction.
\end{proof}

\begin{proposition}\label{prop:positive-density}
Under Definition~\ref{def:GH}, every flow-invariant Borel probability
measure $\mu$ has the form
\[
d\mu=\rho\,d\vol,\qquad \rho\in C^\infty(M),\quad \rho>0.
\]
In particular, every invariant probability measure is equivalent to smooth
volume.
\end{proposition}

\begin{proof}
By Lemma~\ref{lem:invariant-measures-smooth}, $d\mu=\rho\,d\vol$ for a
smooth $\rho\geq0$. Put
\[
Z:=\{x\in M:\rho(x)=0\}.
\]
Equation~\eqref{eq:stationary-density} says that along an orbit,
\[
\frac{d}{dt}\rho(g_t x)
 =-\div_{\vol}X(g_t x)\,\rho(g_t x).
\]
Therefore $\rho(g_tx)=0$ for all $t$ whenever $\rho(x)=0$, so $Z$ is a
compact flow-invariant set. Since $\mu$ is a probability measure,
$Z\neq M$.

Suppose $Z\neq\varnothing$. By the Krylov--Bogolyubov theorem, the compact
flow-invariant set $Z$ supports a flow-invariant probability measure $\eta$.
Lemma~\ref{lem:invariant-measures-smooth} gives
\[
d\eta=\theta\,d\vol,\qquad \theta\in C^\infty(M),\quad\theta\geq0.
\]
Because $\eta\neq0$, there is a point at which $\theta>0$. Continuity then
gives an open subset of $\supp\eta\subset Z$; hence $Z$ has nonempty
interior.

Let
\[
U:=\operatorname{int}Z.
\]
Since $Z$ is invariant, $U$ is invariant. Because $M$ is connected and
$U\neq\varnothing$, while $U\neq M$ since $Z\neq M$, the boundary
$\partial U$ is nonempty. It is compact and flow-invariant. Applying the
Krylov--Bogolyubov theorem once more, choose a flow-invariant probability
measure $\eta_0$ supported on $\partial U$. By
Lemma~\ref{lem:invariant-measures-smooth}, $\eta_0$ is represented by a
nonzero smooth nonnegative function $\theta_0$. Hence
$\supp\eta_0\subset\partial U$ contains a nonempty open subset of $M$.
Thus $\partial U$ has nonempty interior, contradicting
Lemma~\ref{lem:boundary-nowhere}.

Therefore $Z=\varnothing$, which means $\rho>0$.
\end{proof}

\begin{remark}
The same argument shows that the flow is minimal. Indeed, the closure of
every orbit is a compact invariant set and therefore supports an invariant
probability measure. By Proposition~\ref{prop:positive-density}, that
measure has full support, so the orbit closure is $M$. Minimality will not
be needed in the entropy argument below.
\end{remark}

\section{The Greenfield--Wallach spectral theorem}

Let $\mu$ be a flow-invariant probability measure. By
Proposition~\ref{prop:positive-density}, $\mu$ is a smooth positive volume.
Define the Koopman operators
\[
U_t:L^2(M,\mu)\longrightarrow L^2(M,\mu),
\qquad
U_t\varphi:=\varphi\circ g_t.
\]
Since $\mu$ is flow-invariant, $(U_t)_{t\in\mathbb R}$ is a strongly
continuous one-parameter unitary group.

Let
\[
A:=-iX.
\]
For $\varphi,\psi\in C^\infty(M)$, invariance of $\mu$ implies
\[
\int_M (X\varphi)\,\overline{\psi}\,d\mu
 =-\int_M\varphi\,\overline{X\psi}\,d\mu.
\]
Consequently,
\[
\langle A\varphi,\psi\rangle_{L^2(\mu)}
 =\langle\varphi,A\psi\rangle_{L^2(\mu)}.
\]
Thus $A$ is symmetric on the dense domain $C^\infty(M)$. Since $C^\infty(M)$
is a core for the generator of the unitary group $(U_t)$, the closure of
$A$ is self-adjoint and
\[
U_t=e^{itA}.
\]

The crucial analytic input is the following spectral theorem.

\begin{theorem}[Greenfield--Wallach spectral theorem]
\label{thm:GW}
Let $Y$ be a globally hypoelliptic smooth vector field on a compact
connected manifold $M$, and let $\mu$ be a $Y$-invariant probability
measure with smooth positive density. Then the self-adjoint operator
$-iY$ on $L^2(M,\mu)$ has a complete orthonormal system of smooth
eigenfunctions. Equivalently, there exist $\lambda_j\in\mathbb R$ and
$\varphi_j\in C^\infty(M)$ such that
\[
-iY\varphi_j=\lambda_j\varphi_j
\]
and $(\varphi_j)$ is an orthonormal basis of $L^2(M,\mu)$.
\end{theorem}

\begin{proof}
This is the spectral theorem proved by Greenfield and Wallach in
\cite{GreenfieldWallach}. Their global parametrix construction for a
globally hypoelliptic vector field gives a complete spectral resolution by
smooth eigenfunctions for the self-adjoint realization associated with a
smooth positive invariant volume.
\end{proof}

Applying Theorem~\ref{thm:GW} to $Y=X$ yields the following.

\begin{corollary}\label{cor:discrete-spectrum}
For every $g_t$-invariant probability measure $\mu$, the Koopman operator
$U_1$ on $L^2(M,\mu)$ has pure point spectrum.
\end{corollary}

\begin{proof}
By Proposition~\ref{prop:positive-density}, $\mu$ is a smooth positive
volume. Let $(\varphi_j)$ and $(\lambda_j)$ be as in
Theorem~\ref{thm:GW}. Then
\[
U_t\varphi_j=e^{i\lambda_jt}\varphi_j.
\]
In particular,
\[
U_1\varphi_j=e^{i\lambda_j}\varphi_j.
\]
Thus the smooth eigenfunctions $\varphi_j$ form an orthonormal basis of
$L^2(M,\mu)$ consisting of eigenfunctions for $U_1$.
\end{proof}

\section{Discrete spectrum and vanishing entropy}

We use the standard zero-entropy theorem for measure-preserving systems
with discrete spectrum.

\begin{theorem}\label{thm:discrete-zero}
Let $T$ be a measure-preserving transformation of a probability space
$(\Omega,\mathcal B,\mu)$. If the Koopman operator
\[
U_T:L^2(\Omega,\mu)\longrightarrow L^2(\Omega,\mu),\qquad
U_TF=F\circ T
\]
has pure point spectrum, then
\[
h_\mu(T)=0.
\]
\end{theorem}

\begin{proof}
This is the classical discrete-spectrum zero-entropy theorem; see
\cite[Chapter~4]{Walters}. For completeness, in the ergodic case the
Halmos--von Neumann theorem identifies the system measurably with a
translation on a compact metrizable abelian group equipped with Haar
measure. For a left-invariant metric on the group, the translation is an
isometry. An isometry of a compact metric space has zero topological
entropy: the cardinality of an $(n,\varepsilon)$-separated set is
independent of $n$. Hence the corresponding metric entropy is zero. The
general case follows by ergodic decomposition, or directly from the
standard discrete-spectrum entropy theorem.
\end{proof}

\begin{corollary}\label{cor:all-flow-measures-zero}
For every $g_t$-invariant Borel probability measure $\mu$,
\[
h_\mu(g_1)=0.
\]
\end{corollary}

\begin{proof}
By Corollary~\ref{cor:discrete-spectrum}, $U_1$ has pure point spectrum on
$L^2(M,\mu)$. Theorem~\ref{thm:discrete-zero} therefore gives
$h_\mu(g_1)=0$.
\end{proof}

We now use the variational principle in its flow form.

\begin{lemma}\label{lem:flow-variational}
The topological entropy of the time-one map satisfies
\[
\htop(g_1)
 =\sup\bigl\{h_\mu(g_1):
      \mu\text{ is a }g_t\text{-invariant probability measure}\bigr\}.
\]
\end{lemma}

\begin{proof}
The ordinary variational principle \cite[Chapter~8]{Walters} gives the
supremum over all $g_1$-invariant probability measures. If $\mu$ is only
$g_1$-invariant, set
\[
\overline\mu:=\int_0^1 (g_s)_*\mu\,ds.
\]
Since $\mu$ is $g_1$-invariant, $\overline\mu$ is invariant under the whole
flow $(g_t)$. Entropy is invariant under conjugacy and affine in the
invariant measure, so
\[
h_{\overline\mu}(g_1)
 =\int_0^1 h_{(g_s)_*\mu}(g_1)\,ds
 =h_\mu(g_1).
\]
Thus the supremum over $g_1$-invariant measures equals the supremum over
measures invariant under the full flow.
\end{proof}

\begin{theorem}\label{thm:main}
If $X$ is globally hypoelliptic in the sense of
Definition~\ref{def:GH}, then its flow has zero topological entropy:
\[
\htop(g)=\htop(g_1)=0.
\]
\end{theorem}

\begin{proof}
By Corollary~\ref{cor:all-flow-measures-zero},
\[
h_\mu(g_1)=0
\]
for every probability measure $\mu$ invariant under the full flow.
Lemma~\ref{lem:flow-variational} therefore gives
\[
\htop(g_1)=0.
\]
Since the topological entropy of the flow is defined as
$\htop(g_1)$, the proof is complete.
\end{proof}

\begin{remark}
The conclusion is also consistent with the Ruelle--Pesin entropy formula.
Proposition~\ref{prop:positive-density} shows that every invariant measure
is a smooth volume. Once Theorem~\ref{thm:main} gives zero entropy for all
such measures, Pesin's formula implies the vanishing of the sum of the
positive Lyapunov exponents for every invariant measure. A statement of
Pesin's entropy formula for conservative flows can be found in
\cite{EntropyConservativeFlows}. This formula is not used in the proof
above.
\end{remark}

\begin{thebibliography}{9}

\bibitem{GreenfieldWallach}
S.~J. Greenfield and N.~R. Wallach,
\newblock Globally hypoelliptic vector fields,
\newblock \emph{Topology} \textbf{12} (1973), 247--253.

\bibitem{KatokHasselblatt}
A.~Katok and B.~Hasselblatt,
\newblock \emph{Introduction to the Modern Theory of Dynamical Systems},
\newblock Encyclopedia of Mathematics and its Applications \textbf{54},
Cambridge University Press, Cambridge, 1995.

\bibitem{Walters}
P.~Walters,
\newblock \emph{An Introduction to Ergodic Theory},
\newblock Graduate Texts in Mathematics \textbf{79}, Springer-Verlag, New
York, 1982.

\bibitem{EntropyConservativeFlows}
\emph{On the entropy of conservative flows},
\newblock arXiv:1002.2366,
\newblock \url{https://ar5iv.labs.arxiv.org/html/1002.2366}.

\end{thebibliography}

\end{document}